\chapter{Related Work}
\label{sec:related_work}

This thesis sits at the meeting point of three lines of research: Bayesian
experimental design and its amortized variants, simulation-based inference, and the
numerical treatment of stiff differential equations. We review each in enough detail
to make the gap precise. The gap is specific. Amortized policy methods are effective
on standard benchmarks but lose their training signal on stiff, broad-prior systems
with high information per experiment. Per-trajectory methods avoid that failure but
need simulator gradients that stiff systems do not provide reliably. No existing
method uses the current posterior, like the per-trajectory family, while needing
neither a trained policy nor simulator gradients, and while remaining valid when the
posterior is so concentrated that importance weights collapse. We return to this gap
at the end of the chapter.

\section{Bayesian experimental design and the information objective}
\label{sec:rw-bed}

Bayesian experimental design chooses experiments to maximise the expected information
gain (EIG) about unknown parameters~\citep{lindley1956measure,chaloner1995bayesian,ryan2016review}.
For a design $\xi$ and a parameter $\theta$ with prior $p(\theta)$, the EIG is the
mutual information between $\theta$ and the observation $y$,
\begin{equation}
\label{eq:rw-eig}
\mathrm{EIG}(\xi)=I(\theta;y\mid\xi)
=\mathbb{E}_{p(y\mid\xi)}\!\big[H[p(\theta)]-H[p(\theta\mid y,\xi)]\big]
=\mathbb{E}_{p(\theta)p(y\mid\theta,\xi)}\!\left[\log\frac{p(y\mid\theta,\xi)}{p(y\mid\xi)}\right].
\end{equation}
The difficulty is the marginal likelihood $p(y\mid\xi)=\int p(y\mid\theta,\xi)p(\theta)\ud\theta$,
which is intractable for all but the simplest models. Classical estimators use nested
Monte Carlo, drawing an outer sample of $(\theta,y)$ and an inner sample to
approximate the marginal, at a cost that grows with the required
accuracy~\citep{rainforth2018nesting}. In the \emph{sequential} setting the design at
each step depends on the data seen so far, which turns the problem into a control
problem over a growing, history-dependent design space; \citet{rainforth2024modern}
review the modern landscape. The two families that dominate this sequential setting,
and against which PACE is compared, are amortized policies and per-trajectory methods.

\section{Amortized policy-based design}
\label{sec:rw-policy}

\citet{foster2019variational} introduced variational lower bounds on the EIG that can
be optimised by stochastic gradients, and \citet{foster2020unified} unified these into
a single gradient scheme over the prior-contrastive (PCE) family. Deep Adaptive Design
(DAD)~\citep{foster2021deep} made the leap to amortization: rather than optimising a
design at each step, it trains a single policy network $\pi_\phi$ that maps a history
to the next design, by maximising the trajectory mutual information of
Eq.~\eqref{eq:traj-mi}. Because the trajectory marginal is intractable, DAD optimises
the sequential prior-contrastive bound of Eq.~\eqref{eq:traj-spce}, which draws $L$
contrastive parameters from the prior. After training, a design is produced by a
single forward pass, which is what makes the approach attractive for real-time use.
Implicit DAD (iDAD)~\citep{ivanova2021implicit} extended DAD to likelihood-free
models, where the per-step likelihood is unavailable, by using a contrastive bound
estimated from samples, making the method applicable to implicit simulators.

A central fact about all of these methods is that their objective is a contrastive
lower bound on mutual information, and such bounds have a hard ceiling. Any bound
built from $L$ contrastive samples cannot exceed $\log(L+1)$
nats~\citep{mcallester2020formal,paninski2005asymptotic}; the InfoNCE bound that
underlies the contrastive objective makes this explicit, since it is a classification
bound over $L{+}1$ candidates~\citep{oord2018representation,poole2019variational}. This
ceiling is the starting point of our saturation analysis in
Chapter~\ref{sec:methodology}: it is a property of the prior-contrastive objective
itself, not of any particular network or optimiser, so every method in this family
inherits it.

The family is large and active. Step-DAD~\citep{hedman2025stepdad} is semi-amortized:
it pretrains a DAD policy and then finetunes it at each deployment step with a few
gradient updates, trading deployment compute for better designs, and the authors
report instability of pure DAD on the censored-likelihood CES model.
\citet{blau2022optimizing} cast sequential design as a reinforcement-learning problem
and optimise the policy with deep RL, while \citet{huan2016sequential} use approximate
dynamic programming for the same objective. \citet{shen2025vsoed} give a variational
sequential design method trained by RL, and \citet{belliardo2024modelaware} apply
model-aware RL to design in quantum metrology. More recent work amortizes inference
and design jointly: ALINE~\citep{huang2025aline} and the decision-focused method of
\citet{huang2024amortized} train a single network for both posterior inference and
data acquisition, and JADAI~\citep{bracher2025jadai} jointly amortizes adaptive design
and Bayesian inference. For dynamical systems specifically,
\citet{strouwen2026deep} specialise deep adaptive design to model-based design of
experiments for ODEs, and \citet{treloar2022deep} use deep RL to design experiments
for biological ODE models, reporting gains over Fisher-information baselines. All of
these train against the prior-contrastive objective and therefore share its saturation
behaviour on high-information systems.

\section{Per-trajectory and posterior-guided design}
\label{sec:rw-pertraj}

The second family does not train a policy. It maintains an estimate of the current
posterior and chooses the next design by optimising an acquisition score against that
posterior, re-fitting as data arrive. \citet{kleinegesse2020bayesian} estimate the EIG
for implicit models with mutual-information neural estimation and optimise the design
directly. \citet{klein2026supercharging} show that several EIG estimators can be built
on modern simulation-based-inference density estimators and optimised by
multiple-restart gradient ascent through the simulator, unifying inference and design;
this is the closest per-trajectory method to PACE, and the key difference is that it
relies on simulator gradients, whereas PACE uses only forward evaluations.
\citet{iollo2025posterior} and \citet{iollo2025contrastive} develop posterior-guided
and contrastive acquisition criteria, \citet{bharti2024optimizing} optimise designs
through an amortized posterior, and \citet{zaballa2025bridging} bridge sequential
design and simulation-based inference. \citet{vanlier2014bayesian} is an earlier
instance of the same idea for systems-biology ODEs, sequentially selecting informative
conditions by active learning. These methods share PACE's use of the current
posterior; they differ in that they either retrain a proposal each step or depend on
gradients through the simulator.

\section{Simulation-based inference}
\label{sec:rw-sbi}

PACE draws its contrastives from an amortized neural posterior estimator, so it builds
directly on simulation-based inference (SBI), which performs Bayesian inference when
the likelihood is intractable but the simulator can be sampled~\citep{cranmer2020frontier,deistler2025practical}.
Three families are standard. Neural posterior estimation (NPE) parameterises the
posterior $q_\phi(\theta\mid x)$ directly with a conditional density estimator:
\citet{papamakarios2016fast} used mixture density networks,
\citet{lueckmann2017flexible} replaced them with masked autoregressive flows in a
sequential scheme, and \citet{greenberg2019automatic} introduced automatic posterior
transformation (NPE-C) with a correction that keeps the target exact across sequential
rounds. Neural likelihood estimation learns a surrogate likelihood and samples the
posterior with MCMC~\citep{papamakarios2019sequential}, and neural ratio estimation
learns the likelihood-to-evidence ratio by classification~\citep{hermans2020likelihood,durkan2020contrastive,miller2022contrastive}.
The density estimators themselves are usually normalizing
flows~\citep{papamakarios2021normalizing}; we use neural spline
flows~\citep{durkan2019neural} because they give fast, exact density evaluation, which
the importance weighting in PACE requires. More recent flow-matching posterior
estimators~\citep{lipman2023flow,wildberger2024flow,carrasco2026pawsterior} train a
continuous-time velocity field~\citep{chen2018neural} and scale to high dimensions, but
require an ODE solve per density evaluation, which would dominate the per-step cost of
an importance-weighted acquisition score. The \texttt{sbi} toolkit packages these
methods~\citep{tejerocantero2020sbi}, the \texttt{sbibm} suite benchmarks
them~\citep{lueckmann2021benchmarking}, and amortized architectures such as
BayesFlow~\citep{radev2023bayesflow} handle variable numbers of observations.

Two points from this literature matter for PACE. First, sequential SBI methods keep
their importance weights well behaved by retraining the proposal across rounds, for
example with truncated proposals~\citep{deistler2022truncated}; PACE deliberately does
not do this and instead accepts the decline in effective sample size, which is why the
analysis of weight collapse in Chapter~\ref{sec:methodology} is necessary. Second, the
required importance-sampling sample size grows like the exponential of the
Kullback-Leibler divergence between target and proposal~\citep{chatterjee2018sample},
which is the formal reason PACE cannot be read as a calibrated value estimator under a
concentrated posterior. Because amortized posteriors can be miscalibrated, we test
ours with the standard diagnostics: trust and coverage analysis~\citep{hermans2022trust},
simulation-based calibration~\citep{talts2020validating}, distance-to-random-point
coverage~\citep{lemos2023tarp}, and Pareto-smoothed importance
sampling~\citep{vehtari2024pareto}.

\section{Differentiating through stiff simulators}
\label{sec:rw-stiff}

Gradient-based design optimisation, and policy training by the sPCE objective,
differentiate through the ODE solver. On stiff systems this is the dominant cost and a
source of instability. The adjoint method incurs a large compilation overhead and can
be numerically delicate, which has motivated dedicated discretise-then-optimise and
reversible-solver constructions~\citep{kim2024stiff,caldana2025neural,pal2023neural,kidger2024efficient},
and vanishing-gradient pathologies are documented for stiff
dynamics~\citep{linot2025vanishing}. This literature is the practical motivation for
the forward-only optimiser in PACE: by scoring designs with forward simulations only,
PACE avoids the adjoint entirely and applies to stiff and black-box simulators where
gradient-based design is unstable or infeasible.

\section{Classical optimal design for ODE models and model discrimination}
\label{sec:rw-classical}

Before amortization, optimal design for ODE models was dominated by local criteria
based on the Fisher information matrix, such as D-optimality, and by methods for
discriminating between competing models. \citet{huan2014simulation} developed
simulation-based optimal design for nonlinear systems.
\citet{flassig2012optimal} designed optimal stimulus experiments to discriminate
between biochemical ODE models using the overlap of their predictive distributions,
\citet{hainy2022optimal} cast model discrimination as a classification problem and
compared optimal against equidistant designs, and \citet{braniff2022nloed} provide a
software package for nonlinear optimal design in systems biology using the
Jensen-Shannon divergence between predictive densities. These criteria reappear inside
PACE: the collapsed-weight limit of the PACE score reduces to a Fisher-information
criterion (Chapter~\ref{sec:methodology}), which connects the method to this classical
line even though it is derived from a contrastive estimator.

\section{Summary and the gap}
\label{sec:rw-gap}

The amortized policy family casts sequential design as a single trajectory objective
and trains a network to optimise it, but that objective is a prior-contrastive bound
with a $\log(L+1)$ ceiling, so its training signal vanishes on high-information
systems. The per-trajectory family avoids the ceiling by using the current posterior,
but it pays for simulator gradients that stiff systems do not provide reliably, or it
retrains a proposal at every step. Simulation-based inference supplies the posterior
estimator that PACE reuses, but its own theory says the importance weights must
degrade as the posterior concentrates. The gap, then, is a sequential design method
that takes the current posterior from a one-time SBI estimator, scores designs with
forward simulations only, and stays valid when the importance weights collapse. PACE
is that method, and the next chapter derives it.
